\section{Results and Evaluation}
\label{sec:results}

This chapter reports outcomes against the three-axis evaluation
plan committed in Section~\ref{sec:vvt}. Bracketed placeholders
mark data still in collection at the time of writing and are
resolved before submission.

\subsection{Evaluation design and suitability of approach}
\label{sec:eval-design}

The suite is assessed against the three Introduction-stated
claims rather than a single global metric: \emph{coverage} (OWASP
breadth and modern-stack diversity), \emph{integrity} (the
per-user HMAC flag layer resists flag copying and LLM-assisted
solving under deliberate attack), and \emph{engineering rigour}
(intended exploit paths survive code change). Three tracks map
onto these claims. \textbf{Automated regression} supplies an
objective solvability ground truth that removes the
solver-skill confound Meinsma et al.~\cite{Meinsma2022} identify
in platform-log evaluation. An \textbf{LLM solvability trial}
tests the integrity claim directly: a pre-registered five-model
panel attempts every challenge under passive and agentic
scaffoldings, with cold-probe and null-prompt baselines, and a
flag is recorded as solved only when the submitted string is
byte-identical to the HMAC-derived value for the trial salt and
test user. A \textbf{human walkthrough} captures the process
evidence absent from solver counts: time-to-flag, hint-ladder
consumption, NASA-TLX workload~\cite{HartStaveland1988}, and
free-text feedback.

Two methodological choices are worth naming. First, the LLM track
was pre-registered end-to-end before scoring runs began: protocol
(\texttt{Evaluation/llm/PROTOCOL.md}), success criterion
(\texttt{PLAN.md} \S5), failure rubric, prompt SHA-256 hashes
(\texttt{PROMPT\_HASHES.txt}), per-flag regexes
(\texttt{flag\_regexes.json}), and expected flag map
(\texttt{expected\_flags.json}) are all committed artefacts.
Pre-registration is unusual in educational CTF evaluation but is
borrowed from the LLM-benchmark literature; with one or two
seeds per cell, post-hoc metric choice would dominate noise, and
the pre-registered artefacts remove that degree of freedom.
Second, the success criterion is not a regex match but an exact
byte match against the HMAC-derived per-user flag, which makes a
methodologically correct but hallucinated flag a categorical
failure rather than a partial pass.

\subsection{Results}
\label{sec:eval-results}

\subsubsection{Automated regression}
\label{sec:eval-automated}

Every challenge ships an end-to-end exploit script under
\texttt{CTFs/e2e/}; cumulative source is 1\,454 lines across
nine scripts plus a 67-line shared \texttt{conftest.py} providing
\texttt{wait\_for\_service} (\texttt{conftest.py:13}) and
\texttt{load\_credentials} (\texttt{conftest.py:30}) helpers.
Table~\ref{tab:e2e-coverage} summarises coverage: 24 of the
suite's 25 flags are asserted through the intended chain, with
CTF8's third flag fixture exercised via a separate contract
assertion. Test discovery reports 60 collected node IDs
(\texttt{e2e/.pytest\_cache/v/cache/nodeids}); the
\texttt{README.md} headline of ``56 tests'' counts only the
named flag-and-contract assertions.

\begin{table}[t]
\centering
\caption{End-to-end exploit script coverage. LoC is source lines;
flags counts per-flag byte-match assertions.}
\label{tab:e2e-coverage}
\small
\begin{tabular}{|c|c|c|c||c|c|c|c|}
\hline
\textbf{CTF} & \textbf{LoC} & \textbf{Tests} & \textbf{Flags}
& \textbf{CTF} & \textbf{LoC} & \textbf{Tests} & \textbf{Flags}\\ \hline
1 & 84  & 2 & 1 & 6 & 166 & 10 & 4 \\ \hline
2 & 133 & 2 & 1 & 7 & 145 & 6  & 1 \\ \hline
3 & 211 & 9 & 3 & 8 & 110 & 4  & 3 \\ \hline
4 & 142 & 3 & 1 & 9 & 338 & 12 & 6 \\ \hline
5 & 125 & 8 & 4 & \textbf{Tot.} & \textbf{1454} & \textbf{56} & \textbf{24} \\ \hline
\end{tabular}
\end{table}

\texttt{CTFs/e2e/run\_all.sh:19} iterates over
\texttt{ctf1}\ldots\texttt{ctf9} under
\texttt{pytest -v --tb=short} and exits non-zero on any failure;
a wrapper
(\texttt{Evaluation/llm/scripts/phase4\_baseline.sh}) extends
this with per-CTF stack-up, trial-salt flag regeneration, and
tear-down. The harness asserts \emph{intended} solvability, not
the absence of unintended paths: a refactor that introduces a
second bypass next to the intended exploit will still pass. This
limitation motivates the workflow audits discussed in
Section~\ref{sec:eval-impl} and the human track in
Section~\ref{sec:eval-human}. [Insert most recent
\texttt{run\_all.sh} result: per-CTF pass/fail, timestamp, host
platform.] The \texttt{lastfailed} cache currently records one
historical failure
(\texttt{ctf3\_exploit.py::test\_encryption\_key\_in\_bundle}),
traced to a Vite bundling change rather than a regression in the
intended exploit path; [insert disposition].

\subsubsection{LLM robustness trial}
\label{sec:eval-llm}

The panel
(\texttt{Evaluation/llm/run\_matrix.py:62}) comprises
Anthropic \texttt{claude-sonnet-4-6} and
\texttt{claude-haiku-4-5}, OpenAI \texttt{gpt-5-mini}, and Google
\texttt{gemini-2.5-pro} and \texttt{gemini-2.5-flash}. Flagship
Opus- and GPT-5-class models are deliberately excluded from the
primary panel: at roughly USD~9 per nine-CTF Opus run, a
multi-seed evaluation including those models would exceed the
dissertation's hard \pounds20 API budget by an order of magnitude
(\texttt{PLAN.md}~\S11). A spot-check phase
(\texttt{run\_matrix.py:114}) preserves the option of one full
GPT-5 agentic pass on CTF1 and CTF5 if budget headroom allows.
Each model runs two scaffoldings, controlling for the variable
Ji et al.~\cite{ji2025} identify as dominant in CTF-LLM
performance: \emph{passive} (single turn, doc pack + login HTML,
one candidate flag) and \emph{agentic} (system and user prompt,
four tools, fifteen-turn cap). The cap bounds the longest
verified intended chain --- six stages in CTF9 --- with headroom
for tool-use overhead.

\begin{table}[t]
\centering
\caption{LLM trial run matrix
(\texttt{run\_matrix.py:92}). \emph{Upper bound} is a gate, not
a measurement.}
\label{tab:llm-matrix}
\small
\begin{tabular}{|p{1.8cm}|p{4.5cm}|c|}
\hline
\textbf{Phase} & \textbf{Purpose} & \textbf{Cells} \\ \hline
Cold probe  & 5\,models $\times$ 9\,CTFs passive, no stack & 45 \\ \hline
Pilot       & Sonnet $\times$ \{1,5,9\} $\times$ \{pass,agent\} & 6 \\ \hline
Primary     & 5 $\times$ 9 $\times$ 2 conditions $\times$ 2 seeds & 180 \\ \hline
Spot check  & GPT-5 $\times$ \{1,5\} agentic (opt.) & 2 \\ \hline
Null prompt & Sonnet $\times$ CTF1 passive, doc pack replaced & 1 \\ \hline
Upper bound & \texttt{e2e/run\_all.sh} against trial stack & gate \\ \hline
\end{tabular}
\end{table}

Sub-coding of failed flags follows the six-code rubric at
\texttt{RUBRIC.md}:
\texttt{methodology-correct, flag-hallucinated} (the
integrity-critical code),
\texttt{methodology-correct, flag-absent},
\texttt{methodology-partial},
\texttt{methodology-wrong},
\texttt{truncated},
\texttt{harness-error}. A 20\% deterministic-seed sample is
double-rated and Cohen's~$\kappa$~\cite{Cohen1960} is reported,
with a pre-registered threshold of $\geq 0.7$.

Ahead of the automated trial, a small manual pilot was run with
GPT-5.3, with the author acting as the model's shell and HTTP
client (no per-turn cap). Prompts live at
\texttt{Evaluation/llm/manual/ctf<n>.md} and reports at
\texttt{manual/GPT-5.3/ctf<n>-results.md}. Across the four CTFs
run to date, GPT-5.3 solved CTF1 (3 turns, cookie-role tamper),
CTF2 (16 turns, bundle-reconstruction plus PoW plus JWT forgery),
CTF4 (19 turns, reconstructed callback sink, admin-bot
\texttt{.concat()} payload, report-history exfil), and two of
three flags on CTF3 (30 turns; flag~3 not reached before manual
termination). These are not the primary trial; they function as
a calibration read --- a flagship model omitted on cost grounds
solves three of four attempted challenges --- and the solved
chains match the \texttt{SOLUTIONS.md} intended path in every
case, providing independent evidence that the doc-pack curation
does not over-disclose.

[Insert primary-trial findings once
\texttt{Evaluation/llm/runs/} is populated: 9$\times$5
per-flag pass-rate grids (passive and agentic) with
Clopper--Pearson 95\%~CIs~\cite{ClopperPearson1934}, the
agentic-minus-passive delta, and the integrity row: cold-probe
byte-match rate (pre-registered gate: zero), and the fraction of
failed agentic runs sub-coded
\texttt{methodology-correct, flag-hallucinated}. The latter is
the project's headline integrity result --- each such run is a
methodology-correct attempt that would have passed under a
regex-only success criterion, and the HMAC layer is what caught
it.]

\subsubsection{Human walkthrough}
\label{sec:eval-human}

The protocol is fixed in \texttt{Evaluation/human/PLAN.md} with
the session script at \texttt{SESSION\_PROTOCOL.md}. Each
90-minute session comprises a pre-survey, a non-trial warm-up,
three CTF attempts (one per tier, Latin-square counterbalanced,
18-min time cap), and a post-survey. A staged hint ladder
(\texttt{SESSION\_PROTOCOL.md:33}) bounds cognitive cost without
invalidating unaided solves: hints are offered at minute~6
(direction), 10 (mechanism), 14 (payload family) and 17:30 (full
disclosure). Outcomes are scored as \texttt{solved\_unaided},
\texttt{solved\_assisted}, \texttt{partial}, or \texttt{timeout}
with a derived 0.0--1.0 score. Per-participant metrics include
time-to-flag, hints-by-level, distinct techniques attempted,
NASA-TLX workload (\texttt{POST\_SURVEY.md:19}), a five-item
Likert instrument on realism, novelty, calibration, coherence
and recommendation, and free-text responses. Target cohort is
$n = 6$--$10$ (\texttt{PLAN.md:12}); the study is treated as a
calibration pass and qualitative triangulation rather than a
powered inferential study, and a Durham CS low-risk ethics form
was completed prior to the first session.

[Insert findings: achieved $n$; per-tier median time-to-flag;
hint-consumption breakdown; NASA-TLX means and SDs; Likert
means; proportion of \texttt{solved\_unaided} per CTF. One
qualitative paragraph per CTF naming the breadcrumb most often
cited in the post-survey free-text as ``the artefact that moved
me forward'', cross-checked against the observation rubric.]

\subsubsection{Difficulty calibration}
\label{sec:eval-difficulty}

Tier labels (Basic/Intermediate/Advanced) were assigned at
design time from a heuristic combination of stack count, flag
count, and intended-chain depth; they were not validated against
an instrument. Calibration is therefore constructed post-hoc,
treating the human and LLM tracks as independent estimators of
an unobserved difficulty. [Insert: per-tier median
time-to-flag and hint consumption alongside per-tier mean LLM
agentic pass rate. A monotone ordering on both axes --- Basic
faster and fewer hints and higher LLM pass; Advanced inverse ---
is the weakest claim the labels pull in the right direction.
Non-monotonicities --- especially CTF7 (labelled
Basic/Intermediate) and CTF4 (Intermediate, with a long
admin-bot loop) --- are the calibration findings of interest,
with any tier reassignments proposed in 1--2 paragraphs.]

\subsection{Implementation issues and resolutions}
\label{sec:eval-impl}

Three issues were material to the evaluation itself. First, the
e2e harness's initial naive
\texttt{time.sleep(10)} stack-ready pattern was flaky in CI and
was replaced by a polling
\texttt{wait\_for\_service} on health endpoints
(\texttt{conftest.py:13}), with a 60\,s deadline and 2\,s
interval. Second, flag personalisation interacts non-trivially
with dynamic content: CTF3's AES-CBC leak test is sensitive to
Vite bundle chunking, and the resulting historical failure
against the \texttt{test\_encryption\_key\_in\_bundle} assertion
is recorded in \texttt{.pytest\_cache/v/cache/lastfailed};
[insert disposition]. Third, the LLM trial budget is
cost-bounded rather than capability-bounded: the top of the
market is priced at roughly USD~9 per Opus nine-CTF run and was
deliberately displaced into an optional spot-check phase rather
than expanded seed coverage. The per-challenge
\texttt{workflow.md} files (Section~\ref{sec:dev-audit}) record
the audit cycles that found and fixed unintended paths during
development; CTF4 accumulated the most extensive audit trail
(eleven recorded findings), and audits of smaller challenges
were comparatively lighter.

\subsection{Strengths, limitations and threats to validity}
\label{sec:eval-limits}

\textbf{Strengths.} (i) Every challenge has an executable
intended-exploit assertion, which is strictly stronger than
natural-language solution guides. (ii) The LLM track is
pre-registered end-to-end, including the prompt hashes, flag
regexes, and success criterion, which removes the main degree
of freedom in post-hoc metric choice. (iii) The integrity claim
is tested adversarially rather than asserted: a methodologically
correct but hallucinated flag is caught by byte-match and
sub-coded, and a non-trivial hallucination rate is the positive
evidence the HMAC layer did useful work. (iv) The three tracks
triangulate: automated regression, a capability-frontier model
panel, and human think-aloud all bear on the same suite.

\textbf{Sample size.} The human cohort is small ($n=6$--$10$)
and single-institution; reported effect sizes are descriptive
rather than inferential. The LLM trial uses one or two seeds per
cell; Clopper--Pearson intervals are wide and the trial is
framed as diagnostic, not confirmatory (\texttt{PLAN.md}~\S10).
\textbf{Panel selection bias.} The primary LLM panel omits the
largest available models on hard budget grounds; pass rates
therefore characterise mid-tier models, and the manual GPT-5.3
pilot is the only flagship evidence, with single-seed, no-cap
conditions making it incomparable with the formal trial except
qualitatively. \textbf{Salt secrecy.} The integrity claim
depends on the deployment salt remaining secret; if a student
obtains it through misconfiguration or a shared-host
filesystem read, the HMAC layer collapses to a deterministic
function of public inputs and the property fails.
\textbf{Unintended-path coverage.} Automated tests assert
intended exploits only; one-author workflow audits are the only
systematic search for unintended solutions, and an undiscovered
bypass in CTF1, CTF2 or CTF7 cannot be ruled out.
\textbf{Self-report bias.} NASA-TLX and Likert items are
subjective and prone to acquiescence, especially when the
artefact's author runs the session; where free-text and
behavioural evidence diverge, the behavioural reading is
preferred but the bias is not removed. \textbf{No control
arm.} The ``learning gain over status quo'' claim is made by
construction against the literature critique of fixed-VM images,
not by a randomised between-groups comparison against the
COMP2211 baseline; such a study is identified as further work.
\textbf{Partial pre-registration.} The human track's protocol
is fixed but the analysis plan was not formally locked.
\textbf{Single-rater sub-coding.} Only the 20\% double-rated LLM
sample addresses this; passes are single-rater.

\subsection{Reflection on software engineering process and project organisation}
\label{sec:eval-se}

The suite was delivered under a three-phase plan
(\texttt{CTFs/WORKFLOW.md}): challenge-by-challenge build, stack
isolation and orchestration (\texttt{docker-compose.yml} per
CTF), and evaluation. Git hygiene was maintained on
\texttt{main} throughout with
tagged per-challenge commits, and the
commit log is a useful secondary record of which audit cycle
fixed which unintended path. The evaluation harness was built
before the LLM and human trials, which forced the executable
success criterion to be agreed in code before it could be
argued about in prose --- an intentional application of the
Dewey--Kolb experiential-learning framing~\cite{Dewey1938,kolb1984}
to the project's own process. The principal organisational risk
was over-scoping: nine challenges plus three evaluation tracks
plus pre-registration is a large surface, and the primary LLM
trial was deferred to the latest feasible window to keep the
integrity result current rather than retrospective. In
retrospect the budget-driven panel choice was the largest
trade-off imposed by the dissertation envelope; if a single
decision were to be revisited, it would be to ring-fence an
additional ~\pounds40 for a second agentic seed on a flagship
model, which would halve the width of the integrity-row
confidence interval.

\subsection{Summary of findings}
\label{sec:eval-summary}

[Write last, once the four results
sub-sections are filled. Lead with the headline on each axis,
state which of the three project claims each axis supports and
which it does not, and link forward to further work in the
Conclusion.] Skeleton:
(a) \emph{Coverage} --- the e2e suite asserts 24 of 25 flags
across all nine challenges via the intended chain, [insert
\texttt{run\_all.sh} outcome]; (b) \emph{Integrity} --- [insert
cold-probe byte-match rate and the agentic-run hallucination
fraction]; (c) \emph{Pedagogical suitability} --- [insert
proportion of unaided solves per tier, NASA-TLX tier means, and
one sentence on whether tier labels held up].
